\documentclass{article}
\usepackage{amsmath}
\usepackage{amssymb}
\usepackage{amsthm}
\usepackage{mathabx}
\usepackage{graphicx}
\usepackage{parskip}
\usepackage{mathtools}
\usepackage{enumitem}
\usepackage{listings}
\usepackage{forest}
\usepackage{xcolor}
\usepackage{tikz}
\usepackage[utf8]{inputenc}
\usepackage[bulgarian]{babel}
\graphicspath{ {./images/} }

\begin{document}

\definecolor{mGreen}{rgb}{0,0.6,0}
\definecolor{mGray}{rgb}{0.5,0.5,0.5}
\definecolor{mPurple}{rgb}{0.58,0,0.82}
\definecolor{backgroundColour}{rgb}{0.95,0.95,0.92}

\lstdefinestyle{CStyle}{
    backgroundcolor=\color{backgroundColour}, commentstyle=\color{mGreen}, keywordstyle=\color{magenta},
    numberstyle=\tiny\color{mGray}, stringstyle=\color{mPurple}, basicstyle=\ttfamily ,breakatwhitespace=false, 
    breaklines=true, captionpos=b, keepspaces=true, numbers=left, numbersep=5pt, showspaces=false, showstringspaces=false,
    showtabs=false, tabsize=2, language=C, texcl=true, mathescape=true
}

\newcommand{\cplain}[1]{\textnormal{\textcolor{black}{\texttt{#1}}}}

\renewcommand{\proofname}{\textbf{Доказателство:}\vspace{2em}}
\newcommand*\circled[1]{\tikz[baseline=(C.base)]\node[draw,circle,inner sep=1.2pt,line width=0.2mm,](C) {\small #1};\!}

\definecolor{lightBlue}{HTML}{1E90FF}

\title{Теми за Държавен Изпит}
\author{Д. Иванов, КН}
\maketitle

\section*{\centering \color{lightBlue}{Дискретни структури}}

\subsection*{\color{lightBlue}{1. Множества. Декартово произведение. Релации. Функции.}}

\textbf{\underline{Аксиоми} (Цермело-Френкел)}
\begin{itemize}
    \item \textbf{за обема}: $\forall x (x \in A \iff x \in B) \implies A = B$
    \item \textbf{за отделянето}: $X' = \{x \hspace{0.1cm} | \hspace{0.1cm} x \in X, \pi(x)\}$ е множество, където $\pi(x)$ е свойство, чрез което се отделят елементи от $X$
    \item \textbf{за степенното множество}: $\forall X \hspace{0.1cm} \exists \mathcal{P}(X) : \mathcal{P} = \{A \hspace{0.1cm} | \hspace{0.1cm} A \subseteq X\}$
    \item \textbf{за индукцията}: $X = <X_0, \mathcal{F}>$ е множество, където $X_0 \neq \varnothing$, а $\mathcal{F}$ е
    множество от операции, чрез които се изграждат нови елементи на $X$
\end{itemize}
\textbf{\underline{Математическа индукция}} \newline
Метод за доказване на твърдения за индуктивно дефинирани множества. Схема на доказателство по индукция: \newline\newline
\textbf{Теорема}\newline
За всеки елемент $x$ на индуктивно дефинирано множество $X$ е в сила $\pi(x)$
    \begin{itemize}
        \item База: За всеки базов елемент $x_0 \in X_0$ проверяваме верността на $\pi(x_0)$.
        \item Индукционно предположение: Допускаме, че $\pi(x)$ е в сила за всеки елемент $x$, включен до определен
        момент в множеството $X$.
        \item Индукционна стъпка: Показваме, че при направеното предположение $\pi(y)$ също е в сила, $\forall y \in X$
        построен с помощта на операциите от $\mathcal{F}$.
        \item Заключение: $\pi$ е в сила за всеки елемент на множеството $X$.
    \end{itemize}
\textbf{Основни операции} върху множества:
\begin{itemize}
    \item Обединение: $A \cup B = \{x \hspace{0.1cm} | \hspace{0.1cm} x \in A \text{ или } x \in B\}$
    \item Сечение: $A \cap B = \{x \hspace{0.1cm} | \hspace{0.1cm} x \in A \text{ и } x \in B\}$
    \item Разлика: $A \backslash B = \{x \hspace{0.1cm} | \hspace{0.1cm} x \in A \text{ и } x \notin B\}$
    \item Симетрична разлика: $A \triangle B = \{x | (x \in A \text{ и } x \notin B) \text{ или } (x \notin A \text{ и } x \in B)\}$
    \item Допълнение: $\overline{A} = \{x \hspace{0.1cm} | \hspace{0.1cm} x \in U \text{ и } x \notin A\}$
\end{itemize}
\textbf{Свойства на операциите} върху множества:
\begin{itemize}
    \item Идемпотентност: $A \cup A = A, A \cap A = A$
    \item Комутативност: $A \cup B = B \cup A, A \cap B = B \cap A, A \triangle B = B \triangle A$
    \item Асоциативност: $(A \cup B) \cup C = A \cup (B \cup C), (A \cap B) \cap C = A \cap (B \cap C), (A \triangle B)
    \triangle C = A \triangle (B \triangle C)$
    \item Дистрибутивност: $(A \cup B) \cap C = (A \cap C) \cup (B \cap C), (A \cap B) \cup C = (A \cup C) \cap (B \cup C)$
    \item Свойства на празното и универсалното множество: $A \cup \varnothing = A, A \cap \varnothing = \varnothing, A \cup U = U, A \cap U = A$
    \item Свойства на допълнението: $A \cup \overline{A} = U, A \cap \overline{A} = \varnothing, \overline{\overline{A}} = A$
    \item Закони на Де Морган: $\overline{A \cup B} = \overline{A} \cap \overline{B}, \overline{A \cap B} = \overline{A} \cup \overline{B}$
\end{itemize}
\textbf{\underline{Дефиниция}} Нека $a$ и $b$ са произволни елементи. Множеството $(a, b) = \{a, \{a, b\}\}$ наричаме \textbf{наредена двойка} от елементите
$a$ и $b$. \newline\newline
\textbf{\underline{Дефиниция}} Нека $A$ и $B$ са множества. Множеството $A \bigtimes B = \{(a, b) \hspace{0.1cm} | \hspace{0.1cm} a \in A, b \in B\}$
наричаме \textbf{Декартово произведение} на множествата $A$ и $B$. \newline\newline
\textbf{\underline{Дефиниция}}
Нека $A_1,...,A_n$ са множества, $n \in \mathbb{N}$. \textbf{Наредена $n$-орка} наричаме подредена последователност $(a_1, ..., a_n)$, където $a_i \in A_i,
\forall i = 1,...,n$. Множеството $A_1 \bigtimes ... \bigtimes A_n = \{(a_1, ..., a_n) \hspace{0.1cm} | \hspace{0.1cm} a_i \in A_i \hspace{0.1cm} \forall
i = 1,...,n\}$ наричаме \textbf{обобщено Декартово произведение} на множествата $A_1, ..., A_n$. \newline\newline
\textbf{\underline{Дефиниция}}
Нека $A_1,...,A_n$ са множества, $n \in \mathbb{N}$. Всяко $R \subseteq A_1 \bigtimes ... \bigtimes A_n$ наричаме
\textbf{$n$-местна релация} над декартовото произведение $A_1 \bigtimes ... \bigtimes A_n$.
При $n = 2$, релацията наричаме двуместна или просто \textbf{релация}. \newline\newline
\textbf{\underline{Дефиниция}}
Нека $R \subseteq A \bigtimes A$ е релация. Казваме, че $R$ е:
\begin{itemize}
    \item \textbf{Рефлексивна} - ако $\forall a \in A, (a, a) \in R$
    \item \textbf{Симетрична} - ако $\forall a, b \in A, a \neq b, (a, b) \in R \implies (b, a) \in R$
    \item \textbf{Транзитивна} - ако $\forall a, b, c \in A, (a, b) \in R \hspace{0.1cm} \& \hspace{0.1cm} (b, c) \in R \implies (a, c) \in R$
    \item \textbf{Антирефлексивна} - ако $\forall a \in A, (a, a) \notin R$
    \item \textbf{Антисиметрична} - ако $\forall a, b \in A, a \neq b, (a, b) \in R \implies (b, a) \notin R$
    \item \textbf{Силно антисиметрична} - ако $\forall a, b \in A, a \neq b$, точно едно от $(a, b) \in R$ или $(b, a) \in R$ е в сила
\end{itemize}
\textbf{\underline{Дефиниция}}
Нека $R \subseteq A \bigtimes A$ е релация. Казваме, че $R$ е \textbf{релация на еквивалентност}, ако е рефлексивна, симетрична и транзитивна. \newline\newline
\textbf{\underline{Дефиниция}}
Нека $R \subseteq A \bigtimes A$ е релация на еквивалентност и $a \in A$. Множеството $[a] = \{b \hspace{0.1cm} | \hspace{0.1cm} b \in A, aRb\}$ наричаме
\textbf{клас на еквивалентност} на $a$ спрямо $R$. \newline\newline
\textbf{\underline{Дефиниция}}
Нека $R \subseteq A \bigtimes A$ е релация. Казваме, че $R$ е \textbf{частична наредба}, ако е рефлексивна, антисиметрична и транзитивна. \newline\newline
\textbf{\underline{Диаграми на Хассе}}
\begin{center}
    \includegraphics[scale=0.1]{../images/Inclusion_ordering.svg.png}
\end{center}
\textbf{\underline{Дефиниция}}
Нека $R \subseteq A \bigtimes A$ е релация. Казваме, че $R$ е \textbf{пълна(линейна)} наредба, ако е рефлексивна, силно антисиметрична и транзитивна. \newline\newline
\textbf{\underline{Дефиниция}}
Нека $R \subseteq A \bigtimes A$ е релация. Казваме, че елементът $a \in A$ е \textbf{минимален} в $R$, ако $\nexists b \in A, b \neq a$ такъв, че $bRa$. \newline\newline
\textbf{\underline{Дефиниция}}
Нека $R \subseteq A \bigtimes A$ е релация. Казваме, че елементът $a \in A$ е \textbf{максимален} в $R$, ако $\nexists b \in A, b \neq a$ такъв, че $aRb$. \newline\newline
\textbf{\underline{Алгоритъм} (топологично сортиране)} \newline\newline
\underline{Дадено}: Множество $A = \{a_1, ..., a_n\}$ и частична наредба $R \subseteq A \bigtimes A$. \newline
\underline{Резултат}: Наредена $n$-орка $(a_{i_1}, ..., a_{i_n})$ от различните елементи на $A$, определяща пълна наредба
$R' \subseteq A \bigtimes A, R \subseteq R'$ така, че $\forall m, k \in \{1, 2, ... n\}, m \le k \rightarrow (a_{i_m}, a_{i_k}) \in R'$. \newline
\underline{Процедура}: \begin{enumerate}[label=\protect\circled{\arabic*}]
    \item $T = R$ (текуща релация), $M = A$ (неподредени елементи), $j = 1$
    \item Намираме $a_{i_j}$ - минимален в $M$ за релацията $T$
    \item $T = T \hspace{0.1cm} \backslash \hspace{0.1cm} \{(a_{i_j}, a_l) \hspace{0.1cm} | \hspace{0.1cm} (a_{i_j}, a_l) \in T\}, \hspace{0.1cm} M = M \hspace{0.1cm}
    \backslash \hspace{0.1cm} \{a_{i_j}\}, \hspace{0.1cm} j = j + 1$
    \item Ако $M \neq \varnothing$ преминаваме към $\circled{2}$, в противен случай - край.
\end{enumerate}
\textbf{\underline{Дефиниция}}
Нека $f \subseteq X \bigtimes Y$ е релация. Казваме, че $f$ е \textbf{частична функция}, ако $\forall a \in X \hspace{0.1cm} \exists$ не повече от 
едно $b \in Y$ такова, че $(a, b) \in f$ \newline\newline
\textbf{\underline{Дефиниция}}
Нека $f \subseteq X \bigtimes Y$ е релация. Казваме, че $f$ е \textbf{тотална функция}, ако $\forall a \in X \hspace{0.1cm} \exists$ точно
едно $b \in Y$ такова, че $(a, b) \in f$ \newline\newline
\textbf{\underline{Дефиниция}}
Нека $f$ е функция. Казваме, че $f$ е \textbf{инекция}, ако $f(x_1) \neq f(x_2), \forall x_1 \neq x_2 \in X$. \newline\newline
\textbf{\underline{Дефиниция}}
Нека $f$ е функция. Казваме, че $f$ е \textbf{сюрекция} на $X$ върху $Y$, ако $\forall b \in Y \hspace{0.1cm} \exists a \in X$ такова, че $f(a) = b$. \newline\newline
\textbf{\underline{Дефиниция}}
Нека $f : X \rightarrow Y$  е тотална функция. Казваме, че $f$ е \textbf{биекция}, ако е едновременно инекция и сюрекция. \newline\newline
\textbf{\underline{Дефиниция}}
Казваме, че множество $A$ е \textbf{крайно}, ако $A = \varnothing$ или $\exists n \in \mathbb{N}, n \ge 1$ и биекция $f : A \rightarrow \{1, 2, ..., n\}$. \newline\newline
\textbf{\underline{Дефиниция}}
\textbf{Кардиналност} на множество $A$ наричаме броя на елементите му, за които при $A = \varnothing, |A| = 0$ и при $A \neq \varnothing, |A| = n$. \newline\newline
\textbf{\underline{Дефиниция}}
Казваме, че множеството $A$ е \textbf{изброимо безкрайно}, ако $\exists \text { биекция } f : A \rightarrow \mathbb{N}$. Ако множеството $A$ е крайно или изброимо
безкрайно казваме, че то \textbf{изброимо}. \newline\newline
\textbf{\underline{Принцип на Дирихле}}
Нека $X$ и $Y$ са крайни множества и $|X| > |Y|$. Тогава $\forall$ тотална функция $f : X \rightarrow Y$ $\exists x_1 \neq x_2 \in X$ такива, че $f(x_1) = f(x_2)$.

\subsection*{\color{lightBlue}{2. Основни комбинаторни принципи и конфигурации. Рекурентни уравнения.}}

\textbf{\underline{Теорема} (принцип на Дирихле)} \newline
Нека $A$ е множество с $n$ елемента (предмета), а $B$ е множество с $m$ елемента (чекмеджета) и $n > m$. Както и да поставим
всички предмети в чекмеджетата, поне в едно чекмедже ще има поне два предмета. \newline\newline
\textbf{\underline{Теорема} (принцип на биекцията)} \newline
Нека $A$ и $B$ са крайни множества и $|A| = n, |B| = m$. Съществува биекция $f : A \rightarrow B \iff n = m$ \newline\newline
\textbf{\underline{Теорема} (принцип на събирането)} \newline
Нека $A$ е крайно множество и $\mathcal{R} = \{S_1, ..., S_n\}$ е разбиване на $A$. Тогава $|A| = \sum_{i = 1}^{k}|S_i|$ \newline\newline
\textbf{\underline{Теорема} (принцип на изваждането)} \newline
Нека $A$ е крайно множество и $B,C \subseteq A$, като $B = A \backslash C$. Тогава $|B| = |A| - |C|$ \newline\newline
\textbf{\underline{Теорема} (принцип на умножението)} \newline
Нека $A$ и $B$ са крайни множества и $|A| = n, |B| = m$. Тогава $|A \bigtimes B| = |A| . |B| = nm$ \newline\newline
\textbf{\underline{Теорема} (принцип на делението)} \newline
Нека $A$ и $C$ са крайни множества и $B = A \bigtimes C$, $|C| \neq 0$. Тогава $|A| = \frac{|B|}{|C|}$ \newline\newline
\textbf{\underline{Теорема} (принцип за включване и изключване)} \newline
Нека $A$ е крайно множество и $A_1, ..., A_n \subseteq A$. Тогава:
$$|\bigcup^n_{i=1} A_i| = \sum_{1 \leq i \leq n}|A_i| - \sum_{1 \leq i \leq j \leq n}|A_i \cap A_j| + ... + (-1)^{n-1}|A_1 \cap ... \cap A_n|$$
\textbf{\underline{Комбинаторни конфигурации с наредба и повторение}} \newline\newline
Нека $0 < m \in \mathbb{Z}$ и $\mathcal{K}_{\text{н,п}}(n, m)$ е множеството от всички наредени $m$-орки с повторение, съставени
от $n$ различни елемента. Всеки елемент може да участва произволен брой пъти. Тогава:
$$|\mathcal{K}_{\text{н,п}}(n, m)| = n^m$$
\textbf{\underline{Комбинаторни конфигурации с наредба и без повторение}} \newline\newline
Нека $1 \le m \le n$ и $\mathcal{K}_{\text{н}}(n, m)$ е множеството от наредените $m$-орки без повторение, образувани от $n$
елемента. Според \textbf{принципа за умножение} на първо място имаме $n$ избора, на второ - $n - 1$,
на трето - $n - 2$ и така разсъждавайки индуктивно, получаваме формулата за броя на конфигурациите с наредба и без повторение
$\mathcal{K}_{\text{н}}(n, m) = n(n - 1)...(n - m + 1)$. Такива конфигурации се наричат \textbf{вариации} на $n$ елемента от $m$-ти клас и се означават с $V_n^m$\newline
При $m = n$, комбинаторната конфигурация се нарича \textbf{пермутация} и $\mathcal{K}_{\text{н}}(n, n) = |P_n| = n(n - 1)...2.1 = n!$ \newline\newline
\textbf{\underline{Комбинаторни конфигурации без наредба и без повторение}} \newline\newline
Нека $n > 0, 0 \le m \le n$ и $\mathcal{K}(n, m)$ е множеството от ненаредените $m$-орки, без повторение, образувани от $n$ елемента. Такива конфигурации се наричат \textbf{комбинации} на $n$ елемента от $m$-ти клас и се
означават с $C_n^m$. Използваме \textbf{принципа на делението}, за да намерим броя на комбинациите на $n$ елемента от $m$-ти клас.
Ако вместо ненаредените $m$-орки без повторение разгледаме наредените, то $\mathcal{K}(n, m)$ ще се разшири до
$\mathcal{K'}(n, m) = \mathcal{K_{\text{н}}}(n, m)$. Тъй като всяка ненаредена $m$-орка в $\mathcal{K'}(n, m)$ ще се среща по толкова
начина, по колкото можем да подредим $m$ елемента, то тя ще се среща $m!$ пъти. $$\implies |\mathcal{K}(n, m)| = \frac{\mathcal{K_{\text{н}}}
(n, m)}{m!} = \frac{n(n - 1) ... (n - m + 1)}{m!} = \frac{n!}{m!(n - m)!}$$ Така получения израз означаваме с $\binom{n}{m}$ и
наричаме \textbf{биномен коефициент}. \newline\newline
\textbf{\underline{Комбинаторни конфигурации без наредба и с повторение}} \newline
$...$ \newline\newline
\textbf{\underline{Теорема} (Нютон)}
$$(x + y)^n = \sum_{m = 0}^{n}\binom{n}{m}x^my^{n-m} \hspace{0.3cm} \forall n \in \mathbb{N} \hspace{0.2cm} \forall x, y \in \mathbb{R}$$
\begin{proof}
При повдигане на бинома $x + y$ на $n$-та степен всеки едночлен от вида $x^my^{n - m}$ ще се получи толкова пъти, по колкото начина
можем да изберем $m$ множители измежду $n$-те, на които ще поставим $x$, а на останалите $n - m$ поставяме $y$. Броят на тези начини е $\binom{n}{m}$.
\end{proof}
В комбинаториката едно тъждество може да бъде доказано както по формален начин, така и по комбинаторен начин, чрез преброяване на
елементите на подходящо избрана конфигурация по два различни начина. Тази техника е известна като \textbf{принцип на двукратното
броене}. \newline\newline
\textbf{\underline{Теорема}}
$$\binom{n}{m} = \binom{n - 1}{m} + \binom{n - 1}{m - 1} \hspace{0.3cm} n, m \in \mathbb{N}, n > m$$
\begin{proof}
Множеството от комбинаторни конфигурации $C^n_m$ с образуващо множество $A, |A| = n$, можем да разбием на две подмножества: $K_1$ - конфигурациите
несъдържащи елемента $n$ и $K_2$ - съдържащите го. В $K_1$ са $m$-елементните подмножества на множеството от
първите $n - 1$ елемента на $A$, т.е. $|K_1| = \binom{n - 1}{m}$. Елементите на $K_2$ се получават като към всяко
$(m - 1)$ - елементно подмножество на множеството от първите $n - 1$ елемента на $A$, добавим $n$-тия елемент, т.е. $|K_2| = \binom{n - 1}{m - 1}$.
От \textbf{принципа за събирането} $\implies \binom{n}{m} = |\mathcal{K}(n, m)| = |K_1| + |K_2| = \binom{n - 1}{m} + \binom{n - 1}{m - 1}$.
\end{proof}
\textbf{\underline{Дефиниция}}
Нека имаме редица $\tilde{a} = (a_0, a_1, ...)$. Казваме, че тя удовлетворява \textbf{линейно рекурентно отношение} от ред $r$, ако
$$a_{i+r} = c_1a_{i+r-1} + c_2a_{i+r-2} + ... + c_ra_i, \forall i \in \mathbb{N}$$
където $c_1, ..., c_r$ са константи, $c_r \neq 0$ и първите $r$ члена $a_0, ..., a_{r-1}$ са дадени. \newline\newline
\textbf{\underline{Алгоритъм} (решаване на ЛРУ с константни коефициенти)} \newline\newline
...
\subsection*{\color{lightBlue}{3. Графи. Дървета. Обхождания на графи.}}

Нека $V = \{v_1, ..., v_n\}$ е крайно множество, елементите на което наричаме върхове и $E = \{e_1, ..., e_m\}$
е крайно множество, елементите на което наричаме ребра. \newline\newline
\textbf{\underline{Дефиниция}}
Функцията $f_G : E \rightarrow V \bigtimes V$, съпоставяща на всяко ребро наредена двойка от върхове, наричаме
\textbf{краен ориентиран мултиграф}. \newline\newline
\textbf{\underline{Дефиниция}}
Нека $G(V, E, f_G)$ е краен ориентиран мултиграф. Казваме, че $G$ е \textbf{краен ориентиран граф}, ако функцията $f_G $ е такава, че
$f(e_i) \neq f(e_j), i \neq j$. \newline\newline
\textbf{\underline{Дефиниция}}
Нека $G(V, E)$ е краен ориентиран граф. Казваме, че $G$ е \textbf{краен неориентиран граф} или просто \textbf{граф}, ако релацията
$E \subseteq V \bigtimes V$ е антирефлексивна и симетрична. \newline\newline
\textbf{\underline{Дефиниция}}
Нека $G(V, E)$ е мултиграф. Казваме, че има \textbf{път} от $v_{i_0}$ до $v_{i_l}$ в $G$ ($v_{i_0} \leadsto v_{i_l}$), ако
съществува последователност от върхове $v_{i_0}...v_{i_l}$ такава, че $(v_{i_j}, v_{i_{j + 1}}) \in E$ и $v_{i_{j - 1}} \neq v_{i_{j + 1}}, \forall j$. Числото l наричаме
дължина на пътя. При $v_{i_0} = v_{i_l}$, пътя наричаме \textbf{цикъл}. \newline\newline
\textbf{\underline{Дефиниция}}
Нека $G(V, E)$ е мултиграф. Казваме, че $G$ е \textbf{свързан}, ако $v_i \leadsto v_j \hspace{0.1cm} \forall v_i, v_j \in V$. \newline\newline
\textbf{\underline{Дефиниция}}
Нека $G(V, E)$ е мултиграф. Казваме, че подграфът $G'(V', E')$ е негова \textbf{свързана компонента}, ако:
\begin{itemize}
    \item $G'$ е свързан, $V' \subseteq V, E' \subseteq E$
    \item $G'$ е максимален по включване ($\nexists v \in V \backslash V'$, такъв че $G'\cup \{v\}$ да остане свързан)
\end{itemize}
\textbf{\underline{Дефиниция}}
Свързан граф без цикли наричаме \textbf{дърво}. \newline\newline
\textbf{\underline{Дефиниция}}
Граф $D(\{r\}, \{\})$ с един връх $r$ и без ребра наричаме дърво с корен $r$ (\textbf{кореново дърво}). \newline\newline
\textbf{\underline{Теорема}} \newline
Всяко кореново дърво е дърво.
\begin{proof}
Прилагаме индукция по дефиницията на дърво с корен.
\begin{itemize}
    \item Тривиалното кореново дърво $D(\{r\}, \{\})$ е свързан граф $\implies$ е дърво.
    \item Допускаме, че дървото $D(V, E)$ с корен $r$ е свързан граф без цикли.
    \item Ще докажем, че дървото $D'(V' = V \cup \{w\}, E' = E \cup \{(u, w)\})$ с корен $r$ е свързан граф без цикли. \newline\newline
    Нека $v_i, v_j \in V$. Според ИП $D$ е свързан граф $\implies$ $v_i \leadsto v_j$ в $D \implies$ има такъв път и в $D' (D \subseteq D')$.
    В $D'$ е вярно, че $v_i \leadsto w, \forall v_i \in V$ тъй като към пътя от $v_i$ до $u$ в $D$ можем да добавим
    реброто $(u, w)$. Но понеже $w \notin V$, то $w$ е краен връх, свързан единствено с $u \implies$ има път в $D'$ между
    всеки два върха и $D'$ е свързан граф. \newline\newline Според ИП в $D$ няма цикли, а за да има цикъл в $D'$, той трябва да съдържа
    реброто $(u, w)$. Това не е възможно, тъй като $w$ е краен връх $(|w| = 1)$, участва само в реброто $(u, w)$ и по дефиниция не може да участва в цикъл.
\end{itemize}
\end{proof}
\textbf{\underline{Теорема}} \newline
Нека $D(V, E)$ е дърво с корен $r$. Тогава $|V| = |E| + 1$.
\begin{proof}
    $ $
    \begin{itemize}
        \item За тривиалното кореново дърво $D(V = {r}, E = \{\})$ имаме \newline $|V| = 1, |E| = 0 \implies |V| = |E| + 1$
        \item Нека $D(V, E)$ е кореново дърво и $|V| = |E| + 1$
        \item Присъединяваме $w$ към върха $u \in V$ и получаваме кореновото дърво $D'(V' = V \cup {w}, E' = E \cup {(u, w)})$. От
        $|V'| = |V| + 1$ и $|E'| = |E| + 1 \implies |V'| - |E'| = |V| - |E| \overset{\text{от ИП}}{=} 1$. Тогава $|V'| = |E'| + 1$
    \end{itemize}
\end{proof}
\textbf{\underline{Дефиниция}}
Нека $G(V, E)$ е граф, а $D(V, E'), E' \subseteq E$ е дърво. Тогава $D$ наричаме \textbf{покриващо дърво} на $G$. \newline\newline
\textbf{\underline{Алгоритъм} (обхождане в ширина)} \newline\newline
\underline{Дадено}: Свързан граф $G(V, E)$. \newline
\underline{Резултат}: Покриващо дърво $D(V, E')$ на $G$. \newline
\underline{Процедура}:
\begin{enumerate}[label=\protect\circled{\arabic*}]
    \item Избираме коренът $r$ за начален връх на обхождането, т.е. $L_0 = \{r\}$. Построяваме дървото
    $D_0(V_0 = \{L_0\}, E_0 = \varnothing)$ и $l = 0$.
    \item Ако $V_l = V$, алгоритъмът спира и търсеното покриващо дърво е $D_l$. В противен случай преминаваме към $\circled{3}$.
    \item Нека $D_l(V_l, E_l)$ е дървото, построено след $l$-тата стъпка и $L_l$ са върховете от $l$-то ниво. Образуваме
    следващото ниво $L_{l+1} = \{v \hspace{0.1cm} | \hspace{0.1cm} v \notin V_l, \exists w \in L_l, (w, v) \in E\}$. Построяваме дървото $D_{l+1}
    (V_{l+1} = V_l \cup L_{l+1}, E_{l+1} = E_l \cup \{(w, v) \hspace{0.1cm} | \hspace{0.1cm} (w, v) \in E, w \in L_l, v \in L_{l+1}\})$, с присъединяване
    на всеки $v \in L_{l+1}$ към съответния $w \in L_l$. Нека $l = l + 1$ и се връщаме към $\circled{2}$.
\end{enumerate}
\textbf{\underline{Алгоритъм} (обхождане в дълбочина)} \newline\newline
\underline{Дадено}: Свързан граф $G(V, E)$. \newline
\underline{Резултат}: Покриващо дърво $D(V, E')$ на $G$. \newline
\underline{Процедура}:
\begin{enumerate}[label=\protect\circled{\arabic*}]
    \item Избираме коренът $r$ за начален връх на обхождането, т.е. $L_0 = \{r\}$. Построяваме дървото $D_0(V_0= \{r\}, E_0 = \varnothing)$. \newline
    Нека $t = r, parent(t)$ е неопределен и $i = 0$.
    \item Ако $V_l = V$, алгоритъмът спира и търсеното покриващо дърво е $D_l$. В противен случай преминаваме към $\circled{3}$.
    \item Нека е построено дървото $D_i(V_i, E_i)$. Търсим $v \in V_i$ такова, че $(t, v) \in E$ и
    \begin{itemize}
        \item ако има такъв, построяваме дървото $D_{i+1}(V_{i+1} = \{V_i \cup \{v\}\}, E_{i+1} = \{E_i \cup \{(t, v)\}\}),$
        Сега $parent(v) = t, t = v, i = i + 1$ и се връщаме към $\circled{2}$.
        \item в противен случай:
        \begin{itemize}
            \item ако $t \neq r$, тогава $t = parent(t)$ и се връщаме към $\circled{2}$.
            \item ако $t = r$ - край, търсеното покриващо дърво е $D_i$.
        \end{itemize}
    \end{itemize}
\end{enumerate}
\textbf{\underline{Дефиниция}}
\textbf{Ойлеров път} в свързания мултиграф $G(V, E)$ наричаме път, който минава еднократно през всяко ребро на мултиграфа.
Ако Ойлеровият път има съвпадащи начало и край, тогава той се нарича \textbf{Ойлеров цикъл}. \newline\newline
\textbf{\underline{Теорема} (Ойлер)} \newline
Свързаният мултиграф $G(V, E)$ е Ойлеров $\iff$ всеки връх на $G$ е с четна степен.
\begin{proof}
    $ $ \newline  
    $\Rightarrow)$ Нека ребрата на $G(V, E)$ образуват Ойлеров цикъл. Тогава всеки връх има четна степен, защото при обхождане на мултиграфа
    по Ойлеровия цикъл, на всяко ребро "влизащо" във връх $v_i$ съответства ребро "излизащо" от $v_i$, а цикълът съдържа всички ребра точно по
    един път. \newline\newline
    $\Leftarrow)$ Нека мултиграфът $G(V, E)$ е свързан и такъв, че всеки връх има четна степен. Ще построим Ойлеров цикъл от всички ребра на
    мултиграфа.
    \begin{enumerate}[label=\protect\circled{\arabic*}]
        \item Нека $v_0$ е произволен връх на $G$, който избираме за начален и текущ.
        \item Докато е възможно, от текущия връх преминаваме към съседен връх $v$ по необходено ребро, което маркираме като
        обходено, а върхът $v$ става новият текущ. Когато текущият връх няма необходени ребра, значи сме се върнали обратно в
        началния връх $v_0$, тъй като всяко влизане във връх има съответно излизане, а степента на всеки връх е четна.
        В противен случай бихме имали противоречие $\implies$ всички обходени на този етап ребра образуват цикъл $C$.
        \item Ако цикълът $C$ съдържа всички ребра - край на обхождането, построен е Ойлеров цикъл.
        \item В противен случай премахваме от мултиграфа ребрата на цикъла $C$. Полученият мултиграф е такъв, че всеки връх е с четна степен.
        Търсим в $C$ връх $v_1$, от който излиза поне едно необходено ребро. Ако допуснем, че такъв връх няма, при условие, че не
        всички ребра са обходени, ще получим противоречие със свързаността на мултиграфа. Повтаряме действията от стъпка $\circled{2}$ с начален връх
        $v_0 = v_1$ и построяваме нов цикъл. Сега разкъсваме двата цикъла в общия им връх $v_1$ и с отъждествяването на получените краища два по два,
        получаваме нов цикъл $C$ и преминаваме към $\circled{3}$.
    \end{enumerate}
    Алгоритъмът ще завърши работа, когато всички ребра се изчерпят и построява Ойлеров цикъл $\implies$ мултиграфът е Ойлеров.
\end{proof}
\textbf{\underline{Теорема} (Ойлеров път)} \newline
Свързаният мултиграф $G(V, E)$ съдържа Ойлеров път, който не е Ойлеров цикъл $\iff$ има точно $2$ върха с нечетна степен.

\section*{\centering \color{lightBlue}{Езици, автомати и изчислимост}}

\subsection*{\color{lightBlue}{4. Kрайни автомати. Регулярни езици.}}
% \subsection*{\color{lightBlue}{4. Характеризация на регулярните езици. Теорема на Майхил-Нероуд.}}
\textbf{\underline{Дефиниция}}
\textbf{Краен детерминиран автомат} наричаме петорката \newline $\mathcal{A} = <Q, \Sigma, q_0, \delta, F>$, където:
\begin{itemize}
    \item $Q$ е крайно множество от състояния
    \item $\Sigma$ е крайна входна азбука
    \item $q_0 \in Q$ е начално състояние
    \item $\delta : Q \bigtimes \Sigma \rightarrow Q$ е частична функция на преходите, пресмятаща следващото състояние
    \item $F \subseteq Q$ е множество от финални състояния
\end{itemize}
Казваме, че $\mathcal{A}$ е \textbf{тотален}, ако $\delta$ е тотална функция. \newline\newline
\textbf{\underline{Дефиниция}}
Нека $\mathcal{A} = <Q, \Sigma, q_0, \delta, F>$ е краен автомат. Казваме, че състоянието $q$ достига състоянието $q'$ чрез
думата $w = a_1...a_n, a_i \in \Sigma$ в $\mathcal{A}$ ($q \xmapsto[\mathcal{A}]{w} q'$), ако $\exists$ път $q, a_1, q_1, ..., a_n, q'$. \newline\newline
\textbf{\underline{Дефиниция}}
Нека $\mathcal{A} = <Q, \Sigma, q_0, \delta, F>$ е краен автомат. Език, разпознаван от автомата $\mathcal{A}$ наричаме
$$L(\mathcal{A}) = \{w \in \Sigma^* \hspace{0.2cm} | \hspace{0.2cm} q_0 \xmapsto[\mathcal{A}]{w} q, \text{ за } q \in F\}$$
\textbf{\underline{Дефиниция}}
Казваме, че $L \subseteq \Sigma^*$ е \textbf{автоматен език}, ако $\exists$ краен автомат $\mathcal{A}$ такъв, че $L = L(\mathcal{A})$. \newline\newline
\textbf{\underline{Теорема}} \newline
Автоматните езици са затворени относно регулярните операции обединение, конкатенация и звезда на Клини. \newline\newline
\textbf{\underline{Дефиниция}}
\textbf{Регулярни езици} над азбука $\Sigma$ дефинираме рекурсивно:
\begin{itemize}
    \item $\varnothing$, $\{\varepsilon\}$ и $\{a\}$ за всяко $a \in \Sigma$ са регулярни езици.
    \item Ако $L_1$ и $L_2$ са регулярни езици, то $L_1 \cup L_2$, $L_1 . L_2$ и $L_1^*$ са също регулярни езици.
\end{itemize}
\textbf{\underline{Теорема} (Клини)} \newline
Всеки автоматен език е регулярен. \newline\newline
\textbf{\underline{Дефиниция}}
Нека $L \subseteq \Sigma^*$. Казваме, че две думи $w, w' \in \Sigma^*$ са еквивалентни по релацията на Майхил-Нероуд над $L$ ($w \sim_L w'$), ако:
$$\forall u \in \Sigma^* (wu \in L \iff w'u \in L)$$
Релацията $\sim_L$ е релация на еквивалентност и е дясно инвариантна. \newline\newline
\textbf{\underline{Дефиниция}}
Всеки език $L \subseteq \Sigma^*$ дефинира фактормножеството $\Sigma^*/\sim_L = \{[w]_L \hspace{0.1cm} | \hspace{0.1cm} w \in \Sigma^*\}$, където
$[w]_L = \{w' \in \Sigma^* \hspace{0.1cm} | \hspace{0.1cm} w \sim_L w'\}$. \newline\newline
\textbf{\underline{Лема}} \newline
Нека $\mathcal{A} = <Q, \Sigma, q_0, \delta, F>$ е тотален КДА, $q \in Q$ и $w, w' \in \Sigma^*$
са такива, че $q_0 \xmapsto[\mathcal{A}]{w} q$ и $q_0 \xmapsto[\mathcal{A}]{w'} q$. Тогава $w \sim_L w'$. \newline\newline
\textbf{\underline{Теорема (Майхил-Нероуд)}} \newline
Езикът $L \subseteq \Sigma^*$ е регулярен $\iff$ множеството $\Sigma^*/\sim_L$ е крайно.
\begin{proof}
$ $ \newline
$\Rightarrow)$ Нека $L \subseteq \Sigma^*$ е регулярен. Тогава $\exists$ ТКДА $\mathcal{A} = <Q, \Sigma, q_0, \delta, F>$ с $L(\mathcal{A}) = L$.
Нека $\mathcal{A}$ има $n$ достижими състояния $q_0, q_1, ..., q_{n-1}$ и да фиксираме $w_0, w_1, ..., w_{n-1}$ такива, че $q_0 \xmapsto[\mathcal{A}]{w_i}
q_i$ за $0 \le i \le n - 1$. Тогава за произволна дума $w \in \Sigma^*$ е в сила $q_0 \xmapsto[\mathcal{A}]{w} q_i$ за някое $0 \le i \le n - 1$. От предната
\textbf{\underline{Лема}} имаме, че $w \in [w_i]$ за някое $0 \le i \le n - 1 \implies \Sigma^*/\sim_L = \{[w_0]_L, [w_1]_L, ..., [w_{n-1}]_L\} \implies
\Sigma^*/\sim_L$ има най-много $n$ елемента $\implies \Sigma^*/\sim_L$ е крайно. \newline\newline
$\Leftarrow)$ Нека $\Sigma^*/\sim_L$ е крайно. Тогава можем да конструираме ТКДА, който разпознава $L$. Дефинираме автомат $\mathcal{A}_L = <Q_L, \Sigma,
q_{0L}, \delta_L, F_L>$ така:
\begin{itemize}
    \item $Q_L = \Sigma^*/\sim_L$
    \item $q_{0L} = [\varepsilon]_L$
    \item $\delta_L = \{[w]_L \xmapsto[\mathcal{A}_L]{a} [wa]_L \hspace{0.1cm} | \hspace{0.1cm} w \in \Sigma^*, a \in \Sigma \}$
    \item $F_L = \{[w]_L \hspace{0.1cm} | \hspace{0.1cm} w \in L\}$
\end{itemize}
Ще докажем, че за произволни $u, w \in \Sigma^*$ е в сила $[u]_L \xmapsto[\mathcal{A}_L]{w} [uw]_L$ \newline
Прилагаме индукция по $|w|$:
\begin{itemize}
    \item Ако $|w| \le 1$, то твърдението е очевидно.
    \item Нека $|w| = n > 1$ и твърдението е вярно за думи с дължина $n - 1$.
    \item Нека $w = w'a$, където $w' \in \Sigma^*$, $a \in \Sigma$. Тогава $|w'| = n - 1 \overset{\text{от ИП}}{\implies} [u]_L \xmapsto[\mathcal{A}_L]{w'} [uw']_L$.
    От $[uw']_L \xmapsto[\mathcal{A}_L]{a} [uw'a]_L \implies [u]_L \xmapsto[\mathcal{A}_L]{w'a} [uw'a]_L$. Оттук, за произволна дума
    $w \in \Sigma^*$ е в сила $q_{0L} \xmapsto[\mathcal{A}_L]{w} [w]_L \implies w \in L(\mathcal{A}_L) \iff [w]_L \in F_L
    \iff w \sim_L w'$ за някое $w' \in L \implies w' \in L \iff w'\varepsilon \in L \iff w\varepsilon \in L \iff w \in L$, откъдето
    $w \in L(\mathcal{A}_L) \iff w \in L$
\end{itemize}
\end{proof}
\subsection*{\color{lightBlue}{5. Контекстносвободни граматики и езици. Стекови автомати.}}
% \subsection*{\color{lightBlue}{5. Лема за разрастването на контекстносвободни езици. Незатвореност на класа на
% контекстносвободните езици относно сечение и допълнение.}}
\textbf{\underline{Дефиниция}}
\textbf{Контекстносвободна граматика} наричаме четворката \newline $G = <\Gamma, \Sigma, R, S>$, където:
\begin{itemize}
    \item $\Gamma$ е азбука на граматиката
    \item $\Sigma$ са терминални символи на граматиката
    \item $R \subseteq (\Gamma \hspace{0.1cm} \backslash \hspace{0.1cm} \Sigma) \bigtimes \Gamma^*$ са правила на граматиката
    \item $S \in \Gamma \hspace{0.1cm} \backslash \hspace{0.1cm} \Sigma$ е начален символ на граматиката
\end{itemize}
\textbf{\underline{Дефиниция}} Нека $G = <\Gamma, \Sigma, R, S>$ е КСГ. Казваме, че думата $u \in \Gamma^*$ се \textbf{извежда} до
думата $u' \in \Gamma^*$ в $G$ ($u \Rightarrow^*_G u'$), ако $\exists u_1, ..., u_{n-1} \in \Gamma^*$ такива, че $u \Rightarrow_G u_1
\Rightarrow_G ... \Rightarrow_G u_{n-1} \Rightarrow_G u'$. При $n = 0$, $u \Rightarrow_G^0 u$. \newline\newline
\textbf{\underline{Дефиниция}}
Нека $G = <\Gamma, \Sigma, R, S>$ е КСГ. \textbf{Дърво на извод} по $G$ дефинираме по рекурсивен начин:
\begin{itemize}
    \item $\forall a \in \Sigma$, $a$ е дърво на извода с корен $a$, височина $0$ и извод $a$
    \item $\forall A \rightarrow_{G} \varepsilon$,
        \raisebox{-0.9em}{\begin{forest}
          for tree={draw=none, align=center, l sep=10pt, s sep=5pt, baseline}
          [A
            [\(\varepsilon\)]
          ]
        \end{forest}} е дърво с корен $A$, височина $1$ и извод $\varepsilon$
    \item $\forall A \rightarrow_{G} X_1...X_n$, където $X_1, ..., X_n$ са символи на G, и всеки избор на дървета $T_1, ..., T_n$
    съответно с корени $X_1, ..., X_n$, височини $h_1, ..., h_n$ и изводи $w_1, ..., w_n$ \begin{center}
        \begin{forest}
          for tree={draw=none, edge={line width=0.4pt}, align=center, anchor=center, parent anchor=south, child anchor=north, inner sep=2pt, l sep=10pt, s sep=5pt,}
          [A
            [\(T_1\)]
            [\(T_2\)]
            [\(...\)]
            [\(T_n\)]
          ]
        \end{forest}
        \end{center} е дърво с корен $A$, височина $1 + max\{h_1, ..., h_n\}$ и извод $w_1, ..., w_n$
\end{itemize}
\textbf{\underline{Дефиниция}}
Казваме, че езикът $L$ е \textbf{контекстносвободен}, ако $\exists$ КСГ $G$, за която $L(G) = \{w \in \Sigma^* \hspace{0.2cm} | \hspace{0.2cm}
S \Rightarrow^*_{G} w\}$ \newline\newline
\textbf{\underline{Теорема}} \newline
Контекстносвободните езици са затворени относно обединение.
\begin{proof}
Нека $L_1$ и $L_2$ са КСЕ над $\Sigma$, $G_1 = <\Gamma_1, \Sigma, R_1, S_1>$ и $G_2 = <\Gamma_2, \Sigma, R_2, S_2>$ са КСГ с $L(G_1) = L_1$,
$L(G_2) = L_2$ и $\Gamma_1 \cap \Gamma_2 = \Sigma$. \newline Разглеждаме КСГ $G = <\Gamma, \Sigma, R, S>$, където:
\begin{itemize}
    \item $\Gamma = \Gamma_1 \cup \Gamma_2 \cup \{S\}$, за $S \notin \Gamma_1 \cup \Gamma_2$
    \item $R = R_1 \cup R_2 \cup \{S \rightarrow S_1 \hspace{0.1cm} | \hspace{0.1cm} S_2\}$
\end{itemize}
\end{proof}
\textbf{\underline{Теорема}} \newline
Контекстносвободните езици са затворени относно конкатенация.
\begin{proof}
Нека $L_1$ и $L_2$ са КСЕ над $\Sigma$, $G_1 = <\Gamma_1, \Sigma, R_1, S_1>$ и $G_2 = <\Gamma_2, \Sigma, R_2, S_2>$ са КСГ с $L(G_1) = L_1$,
$L(G_2) = L_2$ и $\Gamma_1 \cap \Gamma_2 = \Sigma$. \newline Разглеждаме КСГ $G = <\Gamma, \Sigma, R, S>$, където:
\begin{itemize}
    \item $\Gamma = \Gamma_1 \cup \Gamma_2 \cup \{S\}$, за $S \notin \Gamma_1 \cup \Gamma_2$
    \item $R = R_1 \cup R_2 \cup \{S \rightarrow S_1 S_2\}$
\end{itemize}
\end{proof}
\textbf{\underline{Теорема}} \newline
Контекстносвободните езици са затворени относно звезда на Клини.
\begin{proof}
Нека $L$ е КСЕ над $\Sigma$ и $G = <\Gamma, \Sigma, R, S>$ е КСГ с $L(G) = L$. Разглеждаме КСГ $G' = <\Gamma', \Sigma, R', S'>$, където:
\begin{itemize}
    \item $\Gamma' = \Gamma \cup \{S'\}$, за $S' \notin \Gamma$
    \item $R' = R \cup \{S \rightarrow \varepsilon \hspace{0.1cm} | \hspace{0.1cm} S'S\}$
\end{itemize}
\end{proof}
\textbf{\underline{Дефиниция}}
Казваме, че една КСГ е в нормална форма на Чомски, ако за всяко правило $A \rightarrow_{G} \alpha$ е в сила $|\alpha| = 2$. \newline\newline
\textbf{\underline{Лема}} \newline
Нека $G$ е КСГ в НФЧ, $T$ е дърво на извода в $G$ с връх $A$ и извод $w$ и $T'$ е поддърво на $T$ с връх $B$ и извод $\alpha$. Тогава $\exists
w', w'' \in \Sigma^*$ такива, че $w = w'\alpha w''$ и $A \Rightarrow^*_{G} w'Bw''$. Ако $T \notequiv T'$, то $|w'w''| \ge 1$. \newline\newline
\textbf{\underline{Лема} (разрастване на КСЕ)} \newline
Нека $L$ е КСЕ. Тогава $\exists n \in \mathbb{N}^+$ такова, че $\forall w \in L, |w| > n \hspace{0.1cm} \exists x, y, z, u, v \in \Sigma^*$ такива, че
$w = xyzuv$ и
\begin{itemize}
    \item $|yzu| \le n$
    \item $|yu| \ge 1$
    \item $xy^izu^iv \in L, \forall i \ge 0$
\end{itemize}
\begin{proof}
Нека $G$ е КСГ в НФЧ, която генерира думите от $L$ с дължина $\ge$ 2, $m = |\Gamma \hspace{0.1cm} \backslash \hspace{0.1cm} \Sigma|$, да положим $n = 2^m$ и
$w \in L, |w| > n$. Тогава $|w| > 2$ (в $G$ има поне един нетерминал) и значи в $G$ има дърво на извода $T$ с връх
$S$ (начален за $G$) и извод $w$. От това, че $G$ е в НФЧ следва, че извода на дърво с височина $q$, има дължина $\le 2^q$ и от $|w| > 2^m$ следва, че
$h_T > m$. \newline Конструираме редицата от поддървета на $T : T = T_0 \succ T_1 \succ ... \succ T_h$ и нека
$T_i$ е с връх $A_i$, височина $h - i$ и извод $w_i$ за $0 \leq i \leq h$. Редицата $A_{h-m-1}, A_{h-m}, ..., A_{h-1}$ съдържа
$m + 1$ нетерминала и значи поне един от тях се повтаря два пъти (принцип на Дирихле). \newline Нека $h - m - 1 \leq j < k \leq h - 1$ са такива, че $A_j = A_k$. \newline
От $T \succ T_j$ и предната \textbf{\underline{Лема}} следва, че $\exists x, v \in \Sigma^* : w = xw_jv$ и $S \Rightarrow^*_{G} xA_jv$. От $T_j \succ T_k$ и предната
\textbf{\underline{Лема}} следва, че $\exists y, u \in \Sigma^* : w_j = yw_ku$ и $A_j \Rightarrow^*_{G} yA_ku$ като при това $|yu| \ge 1$. Оттук, полагайки
$z = w_k$ и използвайки $A_j = A_k$ получаваме $S \Rightarrow^*_{G} xA_jv \Rightarrow^*_{G} xyA_juv \Rightarrow^*_{G} xyyA_juuv \Rightarrow^*_{G} ...
\Rightarrow^*_{G} xy^iA_ju^iv \Rightarrow^*_{G} xy^izu^iv$. Освен това $yzu = w_j$ и тъй като $h_{T_j} \le m$, то $|yzu| \le 2^m = n$.
\end{proof}
\textbf{\underline{Пример за неконтекстносвободен език}} \newline\newline
    Езикът $L = \{a^nb^nc^n \hspace{0.2cm} | \hspace{0.2cm} n \ge 0\}$ не е КСЕ.
    \begin{proof}
        Нека $w = a^nb^nc^n \in L$. Тогава $|w| = 3n > n$. Нека $x, y, z, u, v \in \Sigma^*$ такива, че $w = xyzuv, |yzu| \le n$ и
        $|yu| \ge 1$. Тогава $yzu$ не съдържа букви $a$ или не съдържа букви $c$ (или и двете). Нека за определеност $yzu$ не
        съдържа букви $c$ (в другия случай се разсъждава аналогично). Да разгледаме $w_0 = xy^0zu^0v = xzv$. Тогава $w_0$ съдържа
        точно $n$ на брой букви $c$, но тъй като $|yu| \ge 1$, то $|w_0| < 3n \implies w_0 \notin L$. От това и \textbf{\underline{Лема}}
        за разрастването следва, че $L$ не е КСЕ.
    \end{proof}
\textbf{\underline{Теорема}} \newline
Контекстносвободните езици не са затворени относно операцията сечение.
\begin{proof}
Нека $L_1 = \{a^kb^kc^s \hspace{0.1cm} | \hspace{0.1cm} k,s \geq 0\}$ и $L_2 = \{a^kb^sc^s \hspace{0.1cm} | \hspace{0.1cm} k,s \geq 0\}$.
$L_1$ е КСЕ, тъй като е конкатенация на КСЕ $\{a^kb^k \hspace{0.1cm} | \hspace{0.1cm} k \geq 0\}$ и РЕ $\{c^s \hspace{0.1cm}
| \hspace{0.1cm} s \geq 0\}$. $L_2$ е КСЕ, тъй като е конкатенация на РЕ $\{a^k \hspace{0.1cm} | \hspace{0.1cm} k \geq 0\}$
и КСЕ $\{b^sc^s \hspace{0.1cm} | \hspace{0.1cm} s \geq 0\}$. Но $L_1 \cap L_2 = \{a^kb^kc^k \hspace{0.1cm} | \hspace{0.1cm} k \geq 0 \}$, който не е КСЕ.
\end{proof}
\textbf{\underline{Теорема}} \newline
Контекстносвободните езици не са затворени относно операцията допълнение.
\begin{proof}
Съгласно правилото на Де Морган, $L_1 \cap L_2 = \overline{{\overline{L_1} \cup \overline{L_2}}}$. От затвореността на КСЕ относно
обединение и това, че не са затворени относно сечение, следва че КСЕ не са затворени относно допълнение.
\end{proof}
\section*{\centering \color{lightBlue}{Дизайн и анализ на алгоритми}}

\subsection*{\color{lightBlue}{6. Сложност на алгоритъм. Асимптотично поведение на целочислени функции 
($O$-, $\Omega$-, $\Theta$-, $o$- и $\omega$-нотация). Сложност на рекурсивни програми.}}

\textbf{\underline{Дефиниция}}
\textbf{Машина на Тюринг} наричаме петорката $M = <Q, \Gamma, q_0, \delta, F>$, където:
\begin{itemize}
    \item $Q$ е крайно множество от състояния
    \item $\Gamma$ е крайна азбука
    \item $q_0 \in Q$ е начално състояние
    \item $\delta : Q \bigtimes \Gamma \rightarrow (Q \cup F) \bigtimes \Gamma \bigtimes \{L, R, S\}$ е функция на преходите
    \item $F$ е множество от финални състояния, $F \cap Q = \varnothing$
\end{itemize}
\textbf{\underline{Дефиниция}}
\textbf{Машина с произволен достъп} наричаме абстрактен изчислителен модел, подобен на машината на Тюринг, но с възможност
за директен достъп до произволна позиция на лентата. Формално може да се представи като наредена шесторка $M = <Q, \Gamma, q_0, \delta, F, D>$,
където $D$ е произволно множество от команди за достъп до клетки от лентата. \newline\newline
\textbf{\underline{Дефиниция}}
\textbf{Език за програмиране} наричаме крайно множество от символи (азбука) и правила за построяване на синтактично коректни програми над тази азбука. \newline\newline
\textbf{\underline{Дефиниция}}
Нека $f_M : A^* \rightarrow A^*$ е тотална функция изчислима с машината на Тюринг $M$. Функцията
$$t_M(n) = \underset{w \in A^*, |w| = n}{max} [\text{ брой стъпки на M при работа върху $w$ }]$$
наричаме \textbf{сложност по време} на $M$ в най-лошия случай, а функцията
$$\tilde{t}_M(n) = \frac{1}{|A^n|}\sum_{w \in A^*, |w| = n} [\text{ брой стъпки на M при работа върху $w$ }]$$
наричаме \textbf{средна сложност по време} на $M$. Функцията
$$s_M(n) = \underset{w \in A^*, |w| = n}{max} [\text{ брой използвани клетки от M при работа върху $w$ }]$$
наричаме \textbf{сложност по памет} на $M$ в най-лошия случай, а функцията
$$\tilde{s}_M(n) = \frac{1}{|A^n|}\sum_{w \in A^*, |w| = n} [\text{ брой използвани клетки от M при работа върху $w$ }]$$
наричаме \textbf{средна сложност по памет} на $M$. \newline\newline
\textbf{\underline{Дефиниция}}
Нека $f : \mathbb{N} \rightarrow \mathbb{R}$. Казваме, че $f$ е асимптотично положителна, ако $\exists n_0 \in \mathbb{N} \hspace{0.1cm}
\forall n \ge n_0 : f(n) > 0$ \newline\newline
\textbf{\underline{Дефиниция}}
$\mathbf{O(g)}$ е множеството от всички функции, които растат асимптотично не по-бързо от $g$. Означение: $f \preceq g, f = O(g), f \in O(g)$.
$$O(g) = \{f \hspace{0.2cm} | \hspace{0.2cm} \exists c > 0 \hspace{0.2cm} \exists n_0 \in \mathbb{N} \hspace{0.2cm} \forall n \ge n_0 : 0 \le f(n) \le c.g(n)\}$$
\textbf{\underline{Дефиниция}}
$\mathbf{\Omega(g)}$ е множеството от всички функции, които растат асимптотично не по-бавно от $g$. Означение: $f \succeq g, f = \Omega(g), f \in \Omega(g)$.
$$\Omega(g) = \{f \hspace{0.2cm} | \hspace{0.2cm} \exists c > 0 \hspace{0.2cm} \exists n_0 \in \mathbb{N} \hspace{0.2cm} \forall n \ge n_0 : 0 \le c.g(n) \le f(n)\}$$
\textbf{\underline{Дефиниция}}
$\mathbf{\Theta(g)}$ е множеството от всички функции, които растат асимптотично еднакво с $g$. Означение: $f \asymp g, f = \Theta(g), f \in \Theta(g)$.
$$\Theta(g) = \{f \hspace{0.2cm} | \hspace{0.2cm} \exists c_1 > 0 \hspace{0.2cm} \exists c_2 > 0 \hspace{0.2cm} \exists n_0
\in \mathbb{N} \hspace{0.2cm} \forall n \ge n_0 : 0 \le c_1.g(n) \le f(n) \le c_2.g(n)\}$$
\textbf{\underline{Дефиниция}}
$\mathbf{o(g)}$ е множеството от всички функции, които растат асимптотично по-бавно от $g$. Означение: $f \prec g, f = o(g), f \in o(g)$.
$$o(g) = \{f \hspace{0.2cm} | \hspace{0.2cm} \forall c > 0 \hspace{0.2cm} \exists n_0 \in \mathbb{N} \hspace{0.2cm} \forall n \ge n_0 : 0 \le f(n) < c.g(n)\}$$
\textbf{\underline{Дефиниция}}
$\mathbf{\omega(g)}$ е множеството от всички функции, които растат асимптотично по-бързо от $g$. Означение: $f \succ g, f = \omega(g), f \in \omega(g)$.
$$\omega(g) = \{f \hspace{0.2cm} | \hspace{0.2cm} \forall c > 0 \hspace{0.2cm} \exists n_0 \in \mathbb{N} \hspace{0.2cm} \forall n \ge n_0 : 0 \le c.g(n) < f(n)\}$$
\textbf{\underline{Свойства}}
\begin{enumerate}
    \item Рефлексивност на $O$, $\Omega$ и $\Theta$: $f \sigma f, \hspace{0.1cm} \forall \sigma \in \{\preceq, \succeq, \asymp\}$
    \item Симетричност на $\Theta$: $f \asymp g \implies g \asymp f$
    \item Антисиметричност на $O$ и $\Omega$: $f \preceq g$ и $f \succeq g \iff f \asymp g$
    \item Транзитивност на $O, \Omega, \Theta, o, \omega$: ако $f \sigma g$ и $g \sigma h$, то $f \sigma h, \hspace{0.1cm} \newline \forall \sigma \in
    \{\preceq, \succeq, \asymp, \prec, \succ\}$
\end{enumerate}
\textbf{\underline{Теореми} (гранични)}
$$\lim_{n \rightarrow \infty}\frac{f(n)}{g(n)} = c > 0 \implies f(n) = \Theta(g(n))$$
$$\lim_{n \rightarrow \infty}\frac{f(n)}{g(n)} = 0 \iff f(n) = o(g(n))$$
$$\lim_{n \rightarrow \infty}\frac{f(n)}{g(n)} = \infty \iff f(n) = \omega(g(n))$$
\textbf{\underline{Теорема} (Мастър)} \newline
Нека $a \ge 1, b > 1, f(n)$ и $T(n)$ са асимптотично положителни функции, $T(n) = a.T(\frac{n}{b}) + f(n), n > 1$. Тогава:
\begin{enumerate}
    \item Ако $f(n) = O(n^{\log_ba - \varepsilon})$ за $\varepsilon > 0$, то $T(n) = \Theta(n^{\log_ba})$
    \item Ако $f(n) = \Theta(n^{\log_ba})$, то $T(n) = \Theta(n^{\log_ba}logn)$
    \item Ако $f(n) = \Omega(n^{\log_ba + \varepsilon})$ за $\varepsilon > 0$ и $\exists c \in (0, 1) \hspace{0.2cm}
    \exists n_0 \in \mathbb{N} \hspace{0.2cm} \forall n \ge n_0 : a.f(\frac{n}{b}) \le c.f(n)$, то $T(n) = \Theta(f(n))$
\end{enumerate}

\subsection*{\color{lightBlue}{7. Алгоритми в графи с тегла на ребрата. Оценки за сложност.}}
Нека $G(V, E)$ е свързан граф и $w : E \rightarrow \mathbb{R}$ е функция с реални стойности, дефинирана по ребрата на $G$. \newline\newline
\textbf{\underline{Дефиниция}}
Нека $D(V, E')$ е покриващо дърво на $G(V, E)$. \textbf{Тегло} на дървото наричаме сумата $$w(D) = \sum_{e \in E'}w(e) \text{ където } w(e) \text{ e тегло на реброто } e$$
\textbf{\underline{Дефиниция}}
Нека $D'(V, E')$ е покриващо дърво на $G(V, E)$. Казваме, че то е \textbf{минимално}, ако $w(D') \le w(D)$ за всяко друго покриващо дърво $D$ на $G$. \newline\newline
\textbf{\underline{Теорема} (за съгласуваното множество)} \newline
Нека $G(V, E)$ е свързан граф с теглова функция на ребрата $w : E \rightarrow \mathbb{R}$, $\varnothing \neq U \subset V$ и $e = (v_i, v_j) \in E$
е такова, че $v_i \in U, v_j \in V \backslash U$ и $e$ има минимално тегло измежду всички такива ребра. Тогава $\exists$ МПД $D(V, E')$
на $G$ такова, че $e \in E'$.
\begin{proof}
Нека $D(V, E')$ е произволно МПД на $G(V, E)$ и да допуснем, че $e \notin E'$. Да разгледаме графа $G'(V, E' \cup \{e\})$. В $G'$ има точно $1$ цикъл
$\implies$ поне едно ребро в този цикъл пресича границата между $U$ и $V \backslash U$ и нека това да бъде реброто $(u, v)$. По условие $e$ има минимално тегло
измежду всички ребра, които пресичат границата $\implies w(e) \leq w(u, v)$. \newline Нека построим дървото $D'(V, E' \cup \{e\} \hspace{0.1cm} \backslash
\hspace{0.1cm}\{(u, v)\})$. $D'$ е покриващо дърво на $G(V, E)$ и $w(D') = w(D) + w(e) - w(u, v)\leq w(D)$. Ако допуснем, че $w(D') < w(D)$,
стигаме до противоречие с това, че $D(V, E')$ е МПД $\implies w(D') = w(D)$.
\end{proof}
\textbf{\underline{Алгоритъм} (Прим)} \newline\newline
\underline{Дадено}: Свързан граф $G(V, E)$ с теглова функция на ребрата $w : E \rightarrow \mathbb{R}$. \newline
\underline{Резултат}: Минимално покриващо дърво $D(V, E')$ на $G$ с корен - зададен връх $r \in V$. \newline
% \underline{Процедура}:
% \begin{enumerate}[label=\protect\circled{\arabic*}]
%     \item Построяваме дървото $D_0(V_0 = \{r\}, E_0 = \varnothing), k = 0$.
%     \item Нека сме построили $D_k(V_k, E_k)$. \newline Търсим реброто $e = (v_i, v_j)$, където $v_i \in V_k, v_j \in V \backslash V_k$
%     с минимално тегло и построяваме $D_{k+1}(V_{k+1} = V_k \cup \{v_j\}, E_{k+1} = E_k \cup \{e\}), k = k + 1$
%     \item Ако $V_k = V$ - край, полученото дърво $D(V, E' = E_k)$ е оптималното. Иначе минаваме към $\circled{2}$.
% \end{enumerate}
\underline{Псевдокод}:
\begin{lstlisting}[style=CStyle]
foreach v $\in$ V
    v.key $\leftarrow \infty$
    $\pi$[v] $\leftarrow$ Nil
s.key $\leftarrow$ 0

Q $\leftarrow$ {} // празна приоритетна min опашка
Q $\leftarrow$ V

while not isempty(Q) do
    x $\leftarrow$ ExtractMin(Q)
    foreach y $\in$ adj[x]
        if y $\in$ Q and y.key > w(x, y)
            y.key $\leftarrow$ w(x, y)
            $\pi$[y] $\leftarrow$ x
return $\pi$[1...n]
\end{lstlisting}
\underline{Оценка на сложността} \newline\newline
% \begin{itemize}
%     \item Време - $O((V + E)logV)$, при реализация с приоритетна опашка (например min-heap)
%     \item Памет - $O(V + E)$, за представяне на графа и помощните структури
% \end{itemize}
Допускаме, че $\cplain{Q}$ е реализирана чрез двоична пирамида. $\cplain{While}$-цикълът се изпълнява
$\Theta\cplain{(n)}$ пъти, но този фактор не влияе на крайната оценка. Вместо да броим колко пъти се извърта $\cplain{while}$-цикъла,
проследяваме изпълнението на $\cplain{for}$-цикъла на редове 11-14. Ред 11 се изпълнява $\Theta\cplain{(n + m)}$ пъти, защото това е просто
обхождане на списъците на съседство. Ред 13 в най-лошия случай се извиква $\Theta\cplain{(n + m)}$ пъти, всяко със сложност $\cplain{lgn}$.
Оттук оценката за сложността е $\Theta\cplain{((n + m)lgn)}$ \newline\newline
\textbf{\underline{Алгоритъм} (Крускал)} \newline\newline
\underline{Дадено}: Свързан граф $G(V, E)$ с теглова функция на ребрата $w : E \rightarrow \mathbb{R}$. \newline
\underline{Резултат}: Минимално покриващо дърво $D(V, E')$ на $G$. \newline
% \underline{Процедура}:
% \begin{enumerate}[label=\protect\circled{\arabic*}]
%     \item Сортираме ребрата на $G$ в нарастващ ред на теглото и нека този ред е $e_1, ..., e_m$
%     \item От всеки връх $v$ на графа образуваме тривиално дърво $D_v(\{v\}, \varnothing)$
%     \item Обхождаме ребрата $e_1, ..., e_m$ в този ред. За всяко ребро $e_i = (v_{i_1}, v_{i_2})$,
%     ако $v_{i_1}$ и $v_{i_2}$ са в различни дървета, съответно $D'(V', E')$ и $D''(V'', E'')$, ги съединяваме в
%     дървото $D(V' \cup V'', E' \cup E'' \cup \{(v_{i_1}, v_{i_2})\})$
% \end{enumerate}
% \underline{Сложност по}:
% \begin{itemize}
%     \item Време - $O(ElogV)$, при използване на сортиране и структура от тип union-find
%     \item Памет - $O(V + E)$, за представяне на графа и помощните структури
% \end{itemize}
\underline{Псевдокод}:
\begin{lstlisting}[style=CStyle]
A $\leftarrow \varnothing$
// направи разбиването $\{\{1\}, \{2\}, ..., \{n\}\}$
// сортирай E по тегла
count $\leftarrow$ 0

while count < n - 1 do
    e = (u, v) // е следващо ребро в сортираната редица
    if component(u) $\neq$ component(v)
        A $\leftarrow$ A $\cup$ {e}
        identify(component(u), component(v))
        count++
return A
\end{lstlisting}
\underline{Оценка на сложността} \newline\newline
За да бъде ефикасен алгоритъмът, трябва проверката $\cplain{component(u)}$ $\neq$ $\cplain{component(v)}$ и операцията $\cplain{identify(component(u),
component(v))}$ да се извършат във време $O\cplain{(lgn)}$. Ако това е постигнато, то $\cplain{while}$-цикълът ще се изпълнява във време
$O\cplain{(mlgn)}$ и алгоритъмът ще има сложност по време $\Theta(\cplain{mlgn})$ заради сортирането. \newline\newline
\textbf{\underline{Дефиниция}}
Претеглена дължина на пътя $\pi = v_{i_0}, ... v_{i_l}$ в $G$ наричаме $$w(\pi) = \sum_{j=0}^{l-1}w(v_{i_j}, v_{i_{j+1}})$$
\textbf{Най-къс път} от $v_{i_0}$ до $v_{i_l}$ наричаме този с най-малка претеглена дължина. \newline\newline
\textbf{\underline{Теорема}} \newline
Нека $G(V, E)$ е свързан граф с теглова функция $w(e) = 1, \forall e \in E$ и $D$ е покриващо дърво на $G$ с корен $v_0$, построено в ширина.
Всеки път от $v_0$ до връх $v$ в $D$ е най-къс път от $v_0$ до $v$ в графа $G$.
\begin{proof}
Ще докажем с индукция по нивата $L_i, i = 0, 1, ..., l$ на построеното в ширина дърво, че дължината на най-късия път от корена $v_0$ до всеки връх
от ниво $L_i$ е точно $i$.
\begin{itemize}
    \item Дължината на най-късия път от $v_0$ до $v_0$ е $0$ и тъй като той е единствен връх в $L_0$, твърдението е в сила.
    \item Нека твърдението е вярно за $L_0, L_1, ..., L_i$, т.е. $\forall u \in L_j, 0 \le j \le i$, най-късият път от $v_0$ до $u$ има дължина $j$.
    \item Ще докажем, че най-късите пътища от $v_0$ до всеки връх от $L_{i+1}$ са с дължина $i + 1$. Да допуснем, че $\exists v \in L_{i+1}$ такъв, че
    $v_0...w, v$ е най-къс път от $v_0$ до $v$ с дължина $k < i + 1$. Тогава $v_0...w$ е най-къс път от $v_0$ до $w$ с дължина $k - 1 < i$. От ИП
    $w \in L_{k-1}$ и алгоритъмът трябваше да постави $v$ в $L_k, k < i + 1$, но това е противоречие с $v \in L_{i+1} \implies$ твърдението е в сила и за $L_{i+1}$.
\end{itemize}
\end{proof}
% Преди да формулираме алгоритмите на Дейкстра и Белман-Форд, ще въведем две помощни процедури, които ще използваме по-нататък: \newline\newline
% ISS(G, s) - инициализира граф, задавайки начални стойности за разстоянията и предшествениците
% \begin{lstlisting}[style=CStyle]
% foreach v $\in$ V(G)
%     v.d $\leftarrow$ $\infty$ // инициализираме разстоянието до върха v като $\infty$
%     v.$\pi$ $\leftarrow$ Nil // посочваме, че върхът v няма предшественик
% s.d $\leftarrow$ 0 // инициализираме разстоянието до началният връх s с 0
% \end{lstlisting}
% Relax(x, y) - извършва релаксация по ребро (x, y), като актуализира разстоянието до y, ако е намерен по-къс път чрез x
% \begin{lstlisting}[style=CStyle]
% if y.d > x.d + w((x, y)) // ако пътят до y през x е по-къс,
% // го актуализира
%     y.d $\leftarrow$ x.d + w((x, y)) // намалява оценката на най-къс
% // път до y
%     y.$\pi$ $\leftarrow$ x // променя предшественика на y към x
% \end{lstlisting}
\textbf{\underline{Алгоритъм} (Дейкстра)} \newline\newline
\underline{Дадено}: Свързан граф $G(V, E)$ с теглова функция на ребрата $w : E \rightarrow \mathbb{R}^+$ и начален връх $v_0 \in V$. \newline
\underline{Резултат}: Дърво на минималните пътища от $v_0$ до останалите върхове в $G$. \newline
% \underline{Процедура}:
% \begin{enumerate}[label=\protect\circled{\arabic*}]
%     \item Разширяваме $c : E \rightarrow \mathbb{R}^+$ до $c^* : V \bigtimes V \rightarrow \mathbb{R}^+$.
%     \item Нека $dist[0] = 0$, $parent[0] = -1$, $U = \{0\}$ и $dist[i] = c^*(0, i)$, $parent[i] = 0, i = 1, 2, ..., n$
%     \item Повтаряме $n - 1$ пъти стъпките:
%     \begin{itemize}
%         \item Избираме връх $j \notin U$ с минимално $dist[j]$ и $U = U \cup \{j\}$
%         \item За всеки $k \notin U$ пресмятаме $dist[k] = min(dist[k], dist[j] + c^*(j, k))$.
%         Ако $min$ е $dist[j] + c^*(j, k)$, тогава $parent[k] = j$
%     \end{itemize}
% \end{enumerate}
\underline{Псевдокод}:
\begin{lstlisting}[style=CStyle]
foreach v $\in$ V
    v.d $\leftarrow$ $\infty$ // инициализираме разстоянието до върха v
    v.$\pi$ $\leftarrow$ Nil // посочваме, че върхът v няма предшественик
s.d $\leftarrow$ 0 // инициализираме разстоянието до началният връх s

S $\leftarrow \varnothing$ // множество от върхове, за които е намерен НКП от $v_0$
Q $\leftarrow$ {} // празна приоритетна min опашка
Q $\leftarrow$ V

while not isempty(Q) do
    x $\leftarrow$ ExtractMin(Q)
    S $\leftarrow$ S $\cup$ {x}
foreach y $\in$ adj[x]
    if y.d > x.d + w(x, y) // ако пътят до y през x е
    // по-къс, го актуализира
        y.d $\leftarrow$ x.d + w(x, y) // намалява оценката на
        // най-къс път до y
        y.$\pi$ $\leftarrow$ x // променя предшественика на y към x
\end{lstlisting}
% \underline{Сложност по}:
% \begin{itemize}
%     \item Време - $O((V + E)logV)$, с приоритетна опашка (например binary heap)
%     \item Памет - $O(V + E)$, за представяне на графа и помощните структури
% \end{itemize}
\underline{Оценка на сложността} \newline\newline
Ако $\cplain{Q}$ е имплементирана с двоична пирамида, то сложността на алгоритъма е като тази при Прим - $\Theta\cplain{((n + m)lgn)}$ \newline\newline
\textbf{\underline{Алгоритъм} (Флойд-Уоршал)} \newline\newline
\underline{Дадено}: Свързан граф $G(V, E)$ с теглова функция на ребрата $w : E \rightarrow \mathbb{R}^+ \cup \{+\infty\}$ \newline
\underline{Резултат}: Матрица $D$, където $D[i][j]$ съдържа теглото на най-късия път от върха $i$ до върха $j$. \newline
% \underline{Процедура}:
% \begin{enumerate}[label=\protect\circled{\arabic*}]
%     \item За всяко ребро $(u, v) \in E$ с тегло $w(u, v)$ задаваме $dist[u][v] = w(u, v)$.
%     \item За всеки връх $v \in V$ задаваме $dist[v][v] = 0$.
%     \item За всяка междинна възлова точка $k \in V$ и всяка двойка върхове $(i, j) \in V \bigtimes V$,
%     ако $dist[i][j] > dist[i][k] + dist[k][j]$, то $dist[i][j] = dist[i][k] + dist[k][j]$, т.е. ако пътят от $i$ до $j$ чрез междинен връх $k$
%     е по-къс от досега известния най-къс път, стойността на $dist[i][j]$ се обновява.
% \end{enumerate}
\underline{Псевдокод}:
\begin{lstlisting}[style=CStyle]
D$^{(0)}$ $\leftarrow$ W
for k $\leftarrow$ 1 to n
    for i $\leftarrow$ 1 to n
        for j $\leftarrow$ 1 to n
            D$^{(k)}$[i, j] $\leftarrow$ min(D$^{(k-1)}$[i, j],
            D$^{(k-1)}$[i, k] + D$^{(k-1)}$[k, j])
return D$^{(n)}$
\end{lstlisting}
\underline{Оценка на сложността} \newline\newline
Сложността по време е очевидно $\Theta\cplain{(n}^3\cplain{)}$.
% \subsection*{\color{lightBlue}{8. Динамично програмиране. Оценки за сложност.}}
% \textbf{Динамичното програмиране} е метод за решаване на задачи чрез комбиниране на решенията на подзадачи. Алгоритъмът решава всяка
% подзадача веднъж и запазва резултата в таблица, избягвайки нуждата от преизчисляване на отговора всеки път, когато решава някоя
% подзадача. \newline\newline
% Алгоритмичната схема "динамично програмиране" може да се прилага в задачи, които имат рекурсивно решение. Ако в дървото на рекурсията има много
% еднакви подзадачи, то тогава би могло да се "запомни" решението на дадена подзадача и следващия път, когато се налжи да се решава тази задача, вместо
% това да се вземе резултатът наготово. Това "запомняне" наричаме \textbf{мемоизация (memoization)}. Важно условие за прилагане на "мемоизация" е
% рекурсивните извиквания да зависят само от параметрите си и всеки път при едни и същи параметри да се получава
% един и същ резултат. Така е сигурно, че колкото и да се решава дадената подзадача, винаги ще се получава резултатът, който е "запомнен" в началото. \newline\newline
% \textbf{\underline{Принцип за оптималност и конструиране на решението на задачата} \\ \underline{от решенията на подзадачите}} \newline\newline
% Първата стъпка при решаването на дадена оптимизационна задача чрез динамично програмиране е да се определи сртуктурата на оптималното решение.
% Ако за намиране на оптимално решение се използват оптимални решения на по-малки задачи, тогава може да се прилага динамично програмиране. Нещо повече -
% ако това условие не е налице, тогава не е възможно прилагането на динамично програмиране, защото принципът е да свеждаме решаването на задачата до
% решаване на подзадачи от същия вид. Търси се оптимално решение, значи и подзадачите трябва да имат оптимални решения. Тази необходимост оптималното
% решение да се получава от оптимални решения на по-малки задачи се нарича \textbf{"принцип за оптималност"}.
\end{document}